\documentclass[10pt,a4paper]{amsart}
\usepackage{amsmath,amssymb,amsthm,graphicx}
% For commutative diagrams we use Xy-pic
\usepackage[all]{xy}

% We use 2cell for 2-commutative diagrams.
\xyoption{2cell}
% Theorem environments.
%
\theoremstyle{plain}
\newtheorem{theorem}[subsection]{Theorem}
\newtheorem{proposition}[subsection]{Proposition}
\newtheorem{lemma}[subsection]{Lemma}

\theoremstyle{definition}
\newtheorem{definition}[subsection]{Definition}
\newtheorem{example}[subsection]{Example}
\newtheorem{exercise}[subsection]{Exercise}
\newtheorem{situation}[subsection]{Situation}

\theoremstyle{remark}
\newtheorem{remark}[subsection]{Remark}
\newtheorem{remarks}[subsection]{Remarks}

\numberwithin{equation}{subsection}

% Macros
%
\def\lim{\mathop{\mathrm{lim}}\nolimits}
\def\colim{\mathop{\mathrm{colim}}\nolimits}
\def\Spec{\mathop{\mathrm{Spec}}}
\def\Hom{\mathop{\mathrm{Hom}}\nolimits}
\def\Ext{\mathop{\mathrm{Ext}}\nolimits}
\def\SheafHom{\mathop{\mathcal{H}\!\mathit{om}}\nolimits}
\def\SheafExt{\mathop{\mathcal{E}\!\mathit{xt}}\nolimits}
\def\Sch{\mathit{Sch}}
\def\Mor{\mathop{\mathrm{Mor}}\nolimits}
\def\Ob{\mathop{\mathrm{Ob}}\nolimits}
\def\Sh{\mathop{\mathit{Sh}}\nolimits}
\def\NL{\mathop{N\!L}\nolimits}
\def\CH{\mathop{\mathrm{CH}}\nolimits}
\def\proetale{{pro\text{-}\acute{e}tale}}
\def\etale{{\acute{e}tale}}
\def\QCoh{\mathit{QCoh}}
\def\Ker{\mathop{\mathrm{Ker}}}
\def\Im{\mathop{\mathrm{Im}}}
\def\Coker{\mathop{\mathrm{Coker}}}
\def\Coim{\mathop{\mathrm{Coim}}}

% Boxtimes
%
\DeclareMathSymbol{\boxtimes}{\mathbin}{AMSa}{"02}

%
% Macros for moduli stacks/spaces
%
\def\QCohstack{\mathcal{QC}\!\mathit{oh}}
\def\Cohstack{\mathcal{C}\!\mathit{oh}}
\def\Spacesstack{\mathcal{S}\!\mathit{paces}}
\def\Quotfunctor{\mathrm{Quot}}
\def\Hilbfunctor{\mathrm{Hilb}}
\def\Curvesstack{\mathcal{C}\!\mathit{urves}}
\def\Polarizedstack{\mathcal{P}\!\mathit{olarized}}
\def\Complexesstack{\mathcal{C}\!\mathit{omplexes}}
% \Pic is the operator that assigns to X its picard group, usage \Pic(X)
% \Picardstack_{X/B} denotes the Picard stack of X over B
% \Picardfunctor_{X/B} denotes the Picard functor of X over B
\def\Pic{\mathop{\mathrm{Pic}}\nolimits}
\def\Picardstack{\mathcal{P}\!\mathit{ic}}
\def\Picardfunctor{\mathrm{Pic}}
\def\Deformationcategory{\mathcal{D}\!\mathit{ef}}
\usepackage{cleveref}
\crefname{originaldefinition}{original definition}{original definitions}
\Crefname{originaldefinition}{Original definition}{Original definitions}
\makeatletter
\newlabel{topology-definition-spectral-space}{{08YG}{source}{}{sourcepointer-1}{}}
\newlabel{topology-definition-spectral-space@cref}{{[originaldefinition][0][]08YG}{[1][0][]source}{}{}{}}
\makeatother
\crefname{originallemma}{original lemma}{original lemmas}
\Crefname{originallemma}{Original lemma}{Original lemmas}
\makeatletter
\newlabel{topology-lemma-constructible-stable-specialization-closed}{{0903}{source}{}{sourcepointer-2}{}}
\newlabel{topology-lemma-constructible-stable-specialization-closed@cref}{{[originallemma][0][]0903}{[1][0][]source}{}{}{}}
\makeatother
\crefname{originallemma}{original lemma}{original lemmas}
\Crefname{originallemma}{Original lemma}{Original lemmas}
\makeatletter
\newlabel{topology-lemma-make-spectral-space}{{0A31}{source}{}{sourcepointer-3}{}}
\newlabel{topology-lemma-make-spectral-space@cref}{{[originallemma][0][]0A31}{[1][0][]source}{}{}{}}
\makeatother
\title{A closed subset with connected intersections with closures splits off a decomposition of a spectral space}
\author{The Stacks Project Authors}
\date{stacks-project:etale-cohomology.tex:local-snapshot-20261011; GNU FDL 1.2 or later}
\begin{document}
\maketitle
\noindent\textit{Editorial note.} All delivered mathematical text is the author's own, copied from the pinned Stacks Project chapter file etale-cohomology.tex; no statement or proof step was rewritten, reordered or removed.\par\medskip
\noindent\textit{Editorial note.} The counted proof is the author's complete argument: the proof is written out in full: the construction of the specialising sets and their constructible closedness, the partition of X from the connectedness hypothesis, and the specialisation argument giving closedness in X.\par\medskip
\noindent\textit{Editorial note.} The label lemma-connected-topological names the counted statement, and the proof environment immediately following it in the source is delivered verbatim as its proof.\par\medskip
\noindent\textit{Editorial note.} The counted statement is carried by the author's own lemma environment in the chapter source.\par\medskip
\noindent\textit{Editorial note.} The statement and proof are copied from the pinned chapter file; the three cited inputs are kept as external pointers with their canonical Stacks tags.\par\medskip
\noindent\textit{Allowed standard prerequisites (editorial).}\begin{itemize}
\item Spectral spaces and the constructible topology. Topology, the definition of a spectral space and the lemma making the set of points specialising to a quasi-compact subset closed in the constructible topology. Supplies the closedness of the sets W\_i in the constructible topology.
\item Specialisation in the constructible topology. Topology, the lemma allowing a specialisation into a subset closed in the constructible topology to be realised by a point of that subset. Supplies the point x\_1 in W\_1 specialising to x, which is what makes W\_1 closed in the original topology.
\end{itemize}\medskip
\ifdefined\hypertarget\hypertarget{sourcepointer-1}{}\fi
\noindent\textit{Source reference.} Original definition 08YG: sources/stacks-chow/tags.txt, source lines 11339--11339. Prerequisite pointer invoked at proof line(s) 15077; it names a standard fact of a sibling chapter that the counted argument uses, and every step of the counted proof built on it is retained verbatim in the delivered body.\par\medskip
\ifdefined\hypertarget\hypertarget{sourcepointer-2}{}\fi
\noindent\textit{Source reference.} Original lemma 0903: sources/stacks-chow/tags.txt, source lines 11396--11396. Prerequisite pointer invoked at proof line(s) 15092; it names a standard fact of a sibling chapter that the counted argument uses, and every step of the counted proof built on it is retained verbatim in the delivered body.\par\medskip
\ifdefined\hypertarget\hypertarget{sourcepointer-3}{}\fi
\noindent\textit{Source reference.} Original lemma 0A31: sources/stacks-chow/tags.txt, source lines 12724--12724. Prerequisite pointer invoked at proof line(s) 15089; it names a standard fact of a sibling chapter that the counted argument uses, and every step of the counted proof built on it is retained verbatim in the delivered body.\par\medskip
\begin{lemma}
\label{lemma-connected-topological}
Let $Z \subset X$ be a closed subset of a topological space $X$.
Assume
\begin{enumerate}
\item $X$ is a spectral space
(Topology, Definition \ref{topology-definition-spectral-space}), and
\item for $x \in X$ the intersection $Z \cap \overline{\{x\}}$
is connected (in particular nonempty).
\end{enumerate}
If $Z = Z_1 \amalg Z_2$ with $Z_i$ closed in $Z$,
then there exists a decomposition $X = X_1 \amalg X_2$ with
$X_i$ closed in $X$ and $Z_i = Z \cap X_i$.
\end{lemma}

\begin{proof}
Observe that $Z_i$ is quasi-compact. Hence the set of points
$W_i$ specializing to $Z_i$ is closed in the constructible topology
by Topology, Lemma \ref{topology-lemma-make-spectral-space}.
Assumption (2) implies that $X = W_1 \amalg W_2$.
Let $x \in \overline{W_1}$. By Topology, Lemma
\ref{topology-lemma-constructible-stable-specialization-closed} part (1)
there exists a specialization $x_1 \leadsto x$ with $x_1 \in W_1$.
Thus $\overline{\{x\}} \subset \overline{\{x_1\}}$ and we see
that $x \in W_1$. In other words, setting $X_i = W_i$ does the job.
\end{proof}

\end{document}
